\section{Results}
\label{sec:results}
We evaluate ranking quality, the value of angular refinement, transfer to a
held collection, and useful operating points under limited document exposure.

\subsection{Broad Retrieval Gains with a Compact Adapter}
\input{generated/dev_macro}
\begin{figure}[t]
  \centering
  \includegraphics[width=\textwidth]{dev_quality.pdf}
  \caption{Development quality at anchor-inclusive paid-document coverage.
  Points average five seeds and then weight the four collection means equally.
  All six adapters and both native references are retained. Lines connect the
  three measured budgets; they do not imply intermediate measurements.}
  \label{fig:dev-quality}
\end{figure}

\method{} achieves the highest macro \ndcg{} among the six evaluated
methods (Table~\ref{tab:dev}). Its gain over Procrustes is $0.0243$, with
family-adjusted interval $[0.0181,0.0308]$. The effect extends across both
migrations: descriptive pair-wise gains are $0.0257$ for MiniLM and $0.0229$
for E5. Each of the four collection-level macros also favors ARC2 over every
baseline. At the finer condition level, ARC2 leads Procrustes in 23 of 24
point means and every other baseline in all 24. Figure~\ref{fig:dev-quality}
shows the complete budget curves.

The improvement concerns top-ranked relevance: macro Recall is $0.8659$
versus Procrustes' $0.8661$, with descriptive difference interval
$[-0.0049,0.0043]$. The single negative Procrustes nDCG point difference is
FiQA/E5 at 75\% ($-0.0036$). Native Full-X and Full-Y provide the
full-corpus operating context.

\subsection{The Second Correction Adds Value without an Inference Stage}
\begin{figure}[t]
  \centering
  \includegraphics[width=\textwidth]{dev_ablation.pdf}
  \caption{Development mechanism controls. Dots show equal-condition macro
  nDCG differences; bars retain the saved 99\% intervals of the separate
  five-contrast family. The last two rows change the penalty at two steps.
  All intervals use the same continuous horizontal scale.}
  \label{fig:ablation}
\end{figure}

A second correction adds $0.0069$ macro nDCG beyond one step, with family
interval $[0.0038,0.0100]$. The point gain appears in 22 of 24 conditions
and in each of the five seed-level macros. Recomputing the tangent target
therefore provides a repeatable development benefit while the final mapper
remains one affine stage. The complete initialized refinement gains $0.0957$
over affine ridge (Fig.~\ref{fig:ablation}); this comparison jointly evaluates
tangent targeting and the retained alignment prior.

The penalty controls favor moderate regularization. Changing $\lambda$ from
10 to 1 gives $+0.0032$, with an interval crossing zero; changing it to 100
gives $-0.0093$, with a wholly negative interval. The evaluated method keeps
$\lambda=10$ and exactly two corrections.

\subsection{Prospective Transfer: Strong E5 Gains and MiniLM Coverage}
\input{generated/trec_means}
\begin{figure}[t]
  \centering
  \includegraphics[width=\textwidth]{trec_quality.pdf}
  \caption{Prospective TREC-COVID at fixed 1,024-anchor supervision:
  (a,b) primary nDCG@10; (c,d) auxiliary Recall@100. Points average five
  seeds. All six adapters and native references are shown, with explicitly
  panel-specific vertical ranges. Recall curves are descriptive.}
  \label{fig:trec-quality}
\end{figure}

On the held collection, E5 ARC2 exceeds the evaluated Converter and affine
ridge at \textbf{all three migration budgets}. At 25\% coverage, the
Converter gain is $0.2572$, with family interval $[0.1757,0.3376]$. At 50\%,
ARC2 reaches $0.5904$ nDCG versus Converter's $0.3590$ and ridge's $0.3281$
(Table~\ref{tab:trec}). Across both migrations, the 30-contrast family supports
five improvements over Converter, four over ridge and two over SimpleLinear.
The table marks all eleven supported comparisons.

MiniLM provides a useful partial-migration operating point: at 75\% coverage,
ARC2 reaches $0.5195$ nDCG versus native Full-Y's $0.5203$, while one quarter
of the gallery still uses proxies. Its Recall@100 point mean is the highest
among the six adapters at each budget, reaching $0.0998$ at 75\% versus
native Full-Y's $0.1052$. E5 reaches $0.0997$ Recall at 75\%, versus Full-Y's
$0.1162$ (Fig.~\ref{fig:trec-quality}). These auxiliary results distinguish
coverage at 100 from relevance near the top; proximity in nDCG is not an
equivalence test.

Geometric alignment remains competitive: before rounding, Procrustes has
higher nDCG point means in all six cells, although its MiniLM/75\% difference
is below $0.0001$. All Procrustes and LA2M nDCG family intervals cross zero.
The confirmed advantages therefore concern the evaluated Converter, ridge
and SimpleLinear configurations.

\subsection{Higher E5 Quality at Lower Document Exposure}
\input{generated/budget_leverage}
A useful budget pattern extends across all five collections
(Table~\ref{tab:budget-leverage}): E5 ARC2 at 25\% coverage exceeds both
Converter and affine ridge at 75\%. On TREC-COVID, the respective nDCG means
are $0.5628$, $0.4098$ and $0.3278$. The advantage is visible even after the
two comparators receive a much larger set of true destination vectors.

This is a post-hoc descriptive comparison of the measured budgets, rather
than an additional member of the matched-budget inferential family. The table
retains 25\% geometric alternatives, 75\% SimpleLinear and both native
references. It identifies observed quality--exposure operating points;
physical savings and the minimum budget needed for equal quality require
separate measurements.

\subsection{Paid Pairs Strengthen MiniLM Adaptation}
\input{generated/allpaid}
Using already-paid pairs improves MiniLM adaptation without increasing the
50\% document budget. ARC2 rises from $0.3075$ to $0.3671$ nDCG, a gain of
$0.0596$ with family interval $[0.0393,0.0806]$. In matched all-paid
comparisons it exceeds Procrustes, LA2M, SimpleLinear and Converter; ridge is
unresolved. Table~\ref{tab:allpaid} makes these family decisions visible.

E5 changes the method ordering: ARC2 trails all-paid Procrustes and LA2M,
exceeds SimpleLinear and ridge, and is unresolved against Converter. Its
fixed-to-all-paid change, $0.3116$ to $0.2919$, is unresolved. Fixed-anchor
MiniLM ARC2 also trails every all-paid baseline in the cross-regime family
comparisons. All-paid fitting retains $\lambda=10$ in the summed objective;
it tests the fixed procedure at larger supervision, rather than isolating
data quantity at a fixed effective penalty.

\paragraph{Evaluate the gallery that will be searched.}
With each fitted map fixed, true-vector replacement raises MiniLM all-paid
nDCG from $0.3510$ to $0.3671$, but lowers E5 from $0.3149$ to $0.2919$.
For E5, Recall falls from $0.6542$ to $0.5595$ and the paid-document share of
top-10 results falls from $0.4916$ to $0.0647$, despite 50\% paid coverage.
These measured replacement effects make mixed-gallery relevance the
appropriate selection criterion: proxy-only quality can favor a different
operating point from the gallery containing native and mapped vectors.
